% ===========================================================================
%  Diagram TikZ untuk rantai penurunan.
%
%  Ditulis untuk audiens yang bukan orang kendali: tiap gambar menerjemahkan
%  SATU persamaan jadi sesuatu yang bisa ditunjuk. Warna mengikuti palet tema
%  (epicbeamer.sty), jadi tidak ada warna baru diperkenalkan di sini.
%
%    \diagfbd     Langkah 1  gaya menyamping hanya dari thrust yang dimiringkan
%    \diagbudget  Langkah 5  dead time + pengereman harus muat di clearance h
%    \diagrow     Langkah 7  hanya komponen MENDEKAT yang dibatasi
%    \diagarch    slide judul  tiga lapis + plant, satu kolom
% ===========================================================================
\usetikzlibrary{arrows.meta, angles, quotes, decorations.pathreplacing, calc}

\tikzset{
  vec/.style    = {-{Stealth[length=2.2mm]}, line width=0.8pt},
  vecfat/.style = {-{Stealth[length=2.8mm]}, line width=1.3pt},
  hair/.style   = {line width=0.4pt},
  lbl/.style    = {font=\scriptsize, inner sep=1.5pt},
  wee/.style    = {font=\tiny, inner sep=1.5pt},
}

% --------------------------------------------------------------- Langkah 1
% Kenapa miring adalah SATU-SATUNYA cara bergerak menyamping.
\newcommand{\diagfbd}{%
\begin{tikzpicture}[x=1cm, y=1cm]
  \def\th{26}\def\L{2.55}
  \pgfmathsetmacro{\tx}{\L*sin(\th)}
  \pgfmathsetmacro{\ty}{\L*cos(\th)}

  \draw[hair, epicpeach, dashed] (0,-1.5) -- (0,3.0);

  % badan quadrotor, dimiringkan \th
  \begin{scope}[rotate=-\th]
    \draw[line width=1.8pt, epicbrick] (-1.0,0) -- (1.0,0);
    \draw[epicclay, line width=1.0pt] (-1.32,0.14) -- (-0.68,0.14);
    \draw[epicclay, line width=1.0pt] ( 0.68,0.14) -- ( 1.32,0.14);
    \fill[epicbrick] (-1.0,0) circle (1.7pt) (1.0,0) circle (1.7pt);
  \end{scope}

  % thrust tegak lurus badan
  \draw[vecfat, epicrust] (0,0) -- (\tx,\ty)
        node[lbl, epicrust, anchor=south west, inner sep=2pt] {$T$};

  % penguraian: dua kaki persegi panjang
  \draw[hair, epicpeach] (\tx,\ty) -- (0,\ty);
  \draw[hair, epicpeach] (\tx,\ty) -- (\tx,0);
  \draw[vec, epicclay] (0,0) -- (0,\ty)
        node[lbl, epicclay, anchor=east, pos=0.68] {$T\cos\theta$};
  \draw[vec, epicclay] (0,\ty) -- (\tx,\ty)
        node[lbl, epicclay, anchor=south, pos=0.5] {$T\sin\theta$};

  \draw[vec, epicbrick] (0,0) -- (0,-1.25)
        node[lbl, epicbrick, anchor=north] {$mg$};

  \coordinate (O) at (0,0);
  \coordinate (V) at (0,0.95);
  \coordinate (B) at ({0.95*sin(\th)},{0.95*cos(\th)});
  % Urutan titik penting: pic 'angle' menyapu BERLAWANAN arah jarum jam, jadi
  % V--O--B akan menggambar 334 derajat (nyaris lingkaran penuh), bukan 26.
  \pic[draw=epicrust, line width=0.6pt, angle radius=8mm,
       "{$\theta$}"{epicrust, font=\scriptsize, inner sep=1pt}] {angle = B--O--V};

  \node[wee, epicclay, anchor=south] at (0,3.0) {vertical};
\end{tikzpicture}}

% --------------------------------------------------------------- Langkah 5
% Anggaran jarak: dead time + pengereman harus muat di dalam clearance h.
\newcommand{\diagbudget}{%
\begin{tikzpicture}[x=1cm, y=1cm]
  \def\Hh{8.0}\def\Td{2.3}

  \fill[epicpeach] (\Hh+0.8,0) circle (0.78);
  \draw[epicclay, line width=0.8pt] (\Hh+0.8,0) circle (0.78);
  \node[wee, epicbrick] at (\Hh+0.8,0) {obstacle};

  % drone
  \draw[line width=1.8pt, epicbrick] (-0.5,0) -- (0.5,0);
  \draw[epicclay, line width=1.0pt] (-0.68,0.13) -- (-0.32,0.13);
  \draw[epicclay, line width=1.0pt] ( 0.32,0.13) -- ( 0.68,0.13);
  \fill[epicbrick] (-0.5,0) circle (1.7pt) (0.5,0) circle (1.7pt);
  \draw[vec, epicrust] (0.85,0) -- (1.85,0)
        node[lbl, epicrust, anchor=south, pos=0.5] {speed $s$};

  % clearance
  \draw[hair, epicclay] (0,-0.9) -- (0,-1.85);
  \draw[hair, epicclay] (\Hh,-0.9) -- (\Hh,-1.85);
  \draw[{Stealth[length=2mm]}-{Stealth[length=2mm]}, epicclay, line width=0.7pt]
        (0,-1.6) -- (\Hh,-1.6);
  \node[lbl, epicclay, fill=white, inner sep=2pt] at (\Hh/2,-1.6)
    {clearance $h$ to the surface};

  % dua ruas anggaran
  \draw[decorate, decoration={brace, amplitude=5pt}, epicrust, line width=0.8pt]
        (0,1.0) -- (\Td,1.0)
        node[midway, lbl, epicrust, anchor=south, yshift=4pt] {$s\,T_d$};
  \draw[decorate, decoration={brace, amplitude=5pt}, epicbrick, line width=0.8pt]
        (\Td,1.0) -- (\Hh,1.0)
        node[midway, lbl, epicbrick, anchor=south, yshift=4pt]
        {$D(s)$, the braking distance};
  \draw[hair, epicpeach, dashed] (\Td,-0.9) -- (\Td,1.0);

  \node[wee, epicrust, anchor=north, align=center] at (\Td/2,-0.35)
    {reacting:\\[-0.3ex] still closing};
  \node[wee, epicbrick, anchor=north] at ({(\Td+\Hh)/2},-0.35) {slowing down};
\end{tikzpicture}}

% --------------------------------------------------------------- Langkah 7
% Satu baris constraint: hanya komponen MENDEKAT yang dibatasi.
\newcommand{\diagrow}{%
\begin{tikzpicture}[x=1cm, y=1cm]
  \fill[epicpeach] (5.3,0) circle (0.8);
  \draw[epicclay, line width=0.8pt] (5.3,0) circle (0.8);
  \node[wee, epicbrick] at (5.3,0) {hazard};

  \fill[epicbrick] (0,0) circle (2.4pt);
  \node[wee, epicbrick, anchor=north east] at (-0.1,-0.12) {agent};

  \draw[hair, epicpeach, dashed] (0,0) -- (4.5,0);
  \draw[vec, epicclay] (0,0) -- (1.35,0);
  \node[lbl, epicclay, anchor=south] at (0.7,0.06) {$\mathbf{n}_{ij}$};

  % kecepatan yang diminta, dan penguraiannya
  \draw[vecfat, epicrust] (0,0) -- (3.0,1.9)
        node[lbl, epicrust, anchor=south west] {$\mathbf{u}_i$};
  \draw[hair, epicbrick, dashed] (3.0,1.9) -- (3.0,0);
  \draw[hair, epicbrick, dashed] (3.0,1.9) -- (0,1.9);
  \draw[vec, epicbrick] (0,0) -- (3.0,0);
  \node[wee, epicbrick, anchor=north] at (2.2,-0.12) {closing component};
  \draw[vec, epicbrick] (0,0) -- (0,1.9);
  \node[wee, epicbrick, anchor=south, rotate=90] at (-0.16,0.95) {tangential};

  % batas
  \draw[epicrust, line width=1.1pt, dash pattern=on 2.5pt off 2pt]
        (3.6,-0.95) -- (3.6,2.5);
  \node[lbl, epicrust, anchor=south, align=center] at (3.6,2.5)
    {capped at $\lambda_{ij}\varphi(h_{ij})$};
  \node[wee, epicclay, anchor=north, align=center] at (2.4,-1.05)
    {only the closing component is capped --- turning is free};
\end{tikzpicture}}

% --------------------------------------------------- kontrak masuk/keluar
% Strip tiga kotak: APA YANG MASUK (dengan satuan) -> APA YANG DIHITUNG ->
% APA YANG KELUAR (dengan satuan). Dipakai sekali di tiap lapis, sehingga
% keluaran satu lapis terlihat persis menjadi masukan lapis berikutnya.
\newcommand{\iobox}[4]{% #1 x  #2 judul  #3 warna judul  #4 isi
  \node[draw=epicboxline, line width=0.5pt, fill=white, rounded corners=1pt,
        minimum width=38mm, minimum height=21mm, align=left, inner sep=2.2mm,
        font=\scriptsize, text=epicbrick, anchor=north west] (b#1) at (#1,0) {#4};
  \node[fill=#3, text=white, font=\bfseries\tiny, inner xsep=2mm, inner ysep=0.9mm,
        anchor=north west] at ($(b#1.north west)+(1.6mm,1.9mm)$) {#2};}

\newcommand{\diagio}[3]{%
\begin{tikzpicture}[x=1mm, y=1mm]
  \iobox{0}{WHAT GOES IN}{epicclay}{#1}
  \iobox{45}{WHAT IS COMPUTED}{epicrust}{#2}
  \iobox{90}{WHAT COMES OUT}{epicbrick}{#3}
  \draw[vec, epicclay] (38.5,-9.5) -- (44,-9.5);
  \draw[vec, epicclay] (83.5,-9.5) -- (89,-9.5);
\end{tikzpicture}}

% ------------------------------------------- kaskade low level + satuannya
% Rantai yang benar-benar dieksekusi 250 kali per detik, dengan SATUAN di tiap
% panah -- itu yang membuat jelas apa yang masuk dan apa yang keluar.
\newcommand{\diagcascade}{%
\begin{tikzpicture}[x=1mm, y=1mm,
   blk/.style={draw=epicrust, line width=0.6pt, fill=epicbeige!35, rounded corners=1pt,
               minimum width=20mm, minimum height=10mm, align=center,
               font=\tiny, text=epicbrick},
   sig/.style={font=\tiny, text=epicclay, inner sep=1pt}]
  \node[blk] (pf) at (0,0)   {pre-filter};
  \node[blk] (op) at (28,0)  {outer PID};
  \node[blk] (ip) at (56,0)  {inner PD};
  \node[blk] (mx) at (84,0)  {mixer};
  \node[draw=epicclay, line width=0.7pt, fill=epicpeach!40, rounded corners=1pt,
        minimum width=20mm, minimum height=10mm, align=center, font=\tiny,
        text=epicbrick] (pl) at (112,0) {quadrotor};

  % jalur maju; label SINYAL di atas panah, bukan di atas kotak
  \draw[vec, epicbrick] (-17,0) -- (pf);
  \node[sig, anchor=south] at (-8.5,6) {$\mathbf{p}^\star$ [m]};
  \node[sig, anchor=north, text=epicrust] at (-8.5,-6) {from mid level};
  \foreach \a/\b/\x/\t in {pf/op/14/{[m], [m/s]}, op/ip/42/{$\theta_{\mathrm{ref}}$ [rad]},
                           ip/mx/70/{$\tau$ [N$\cdot$m]}, mx/pl/98/{$\Omega$ [rad/s]}} {
    \draw[vec, epicbrick] (\a) -- (\b);
    \node[sig, anchor=south] at (\x,6) {\t};}

  % umpan balik: dua lintasan terpisah dari plant
  \draw[vec, epicclay] (pl.south) -- (112,-13) -- (56,-13) -- (ip.south);
  \node[sig, anchor=south] at (84,-12.6) {$\theta,\dot\theta$ [rad], [rad/s]};
  \draw[vec, epicclay] (116,-5) -- (116,-22) -- (28,-22) -- (op.south);
  \draw[epicclay, line width=0.8pt] (pl.south east) ++(-2,0) -- (116,-5);
  \node[sig, anchor=south] at (72,-21.6) {$\mathbf{p},\mathbf{v}$ [m], [m/s]};

  \node[sig, text=epicrust, anchor=north west, align=left] at (-17,-26)
    {the whole chain runs 250 times per second};
\end{tikzpicture}}

% ------------------------------------------------- arsitektur, slide judul
% Gambaran tiga lapis dalam satu kolom: tiap panah ke bawah adalah keluaran
% satu lapis yang menjadi masukan lapis di bawahnya; satu panah balik membawa
% odometri + LiDAR dari simulator ke ketiga lapis.
\newcommand{\diagarch}{%
\begin{tikzpicture}[x=1mm, y=1mm,
   lay/.style={draw=epicboxline, line width=0.5pt, fill=epicboxfill,
               rounded corners=1pt, minimum width=46mm, minimum height=6.4mm,
               align=left, inner xsep=2mm, inner ysep=1.4mm, font=\tiny, text=epicbrick,
               text width=42mm},
   tag/.style={font=\bfseries\tiny, text=white, inner xsep=1.3mm,
               inner ysep=0.6mm, anchor=south west},
   sig/.style={font=\tiny, text=epicclay, inner sep=1pt, anchor=west}]
  \node[lay] (hi) at (0,0)   {Voronoi cells $\to$ boustrophedon lanes\\
                              re-partition when an agent dies};
  \node[lay] (mi) at (0,-12) {CBF-QP safety filter\\
                              obstacles, walls, neighbours};
  \node[lay] (lo) at (0,-24) {\cL{PID-LQR} \textbar{} \cH{PID-$H_\infty$}\\
                              position $\to$ four rotor speeds};
  \node[lay, draw=epicclay, fill=epicpeach!40] (pl) at (0,-36)
       {Gazebo \textperiodcentered{} 7 quadrotors\\
        wind \textperiodcentered{} cylinders \textperiodcentered{} agent loss};
  \foreach \n/\t/\c in {hi/HIGH \textperiodcentered{} event/epicclay,
                        mi/MID \textperiodcentered{} 20 Hz/epicrust,
                        lo/LOW \textperiodcentered{} 250 Hz/epicbrick,
                        pl/PLANT/epicclay}
    \node[tag, fill=\c] at ($(\n.north west)+(1.6mm,-0.9mm)$) {\t};
  \foreach \a/\b/\t in {hi/mi/lane, mi/lo/safe velocity, lo/pl/rotor speeds} {
    \draw[vec, epicbrick] (\a.south) -- (\b.north);
    \node[sig] at ($(\a.south)!0.5!(\b.north)+(1mm,0)$) {\t};}
  % umpan balik ke ketiga lapis
  \draw[vec, epicclay] (pl.east) -- ++(3,0) coordinate (fb) |- (hi.east);
  \foreach \n in {mi,lo} \draw[vec, epicclay] (fb |- \n.east) -- (\n.east);
  \node[font=\tiny, text=epicclay, rotate=90, anchor=north] at ($(fb)+(0.4,18)$)
       {odometry \textperiodcentered{} LiDAR};
\end{tikzpicture}}
